\documentclass[11pt,a4paper]{article}
\usepackage[T1]{fontenc}
\usepackage{lmodern}
\usepackage[margin=25mm,headheight=14pt]{geometry}
\usepackage{amsmath,amssymb,amsthm,mathtools}
\usepackage{booktabs,array,enumitem,microtype,fancyhdr}
\usepackage[hidelinks,pdfusetitle]{hyperref}
\hypersetup{pdftitle={Evaluation of an explicit local-field Haar core: complete formalized statement},pdfauthor={}}
\setlength{\parindent}{0pt}
\setlength{\parskip}{6pt}
\setlist[enumerate]{leftmargin=*,itemsep=7pt,topsep=5pt}
\pagestyle{fancy}
\fancyhf{}
\fancyhead[L]{\small Explicit local integral evaluation}
\fancyhead[R]{\small Complete statement, v1}
\fancyfoot[C]{\thepage}
\renewcommand{\headrulewidth}{0.3pt}
\newcommand{\C}{\mathbb C}
\newcommand{\R}{\mathbb R}
\newcommand{\Z}{\mathbb Z}
\newcommand{\OO}{\mathcal O}
\newcommand{\mm}{\mathfrak m}
\newcommand{\abs}[1]{\lvert#1\rvert}
\newcommand{\Ph}{\Phi_0}
\newcommand{\I}{\mathcal I}
\newcommand{\E}{E_q}
\newtheorem{theorem}{Theorem}
\newtheorem{proposition}{Proposition}
\newtheorem{remark}  [theorem] {Remark}
\title{Evaluation of an explicit local integral}
\author{}
\date{Complete formalized statement\quad\textnormal{Version 1}}
\begin{document}
\thispagestyle{empty}
\begin{center}
{\LARGE\bfseries Evaluation of an explicit local integral}\par\vspace{7pt}
{\large Complete formalized statement}\par\vspace{4pt}
{Version 1}
\end{center}
\vspace{5pt}

This document specifies all data, definitions and hypotheses needed for the
formalized theorem. The function $\Ph$ is defined by an explicit finite formula below.


\section{Field, measures and scalar parameters}

Let $F$ be a nonarchimedean local field. Equip it with its local-field topology
and Borel $\sigma$-algebra, and equip $F^\times$ with the induced topology and
Borel structure. Let
\[
 v:F\longrightarrow\Z\cup\{+\infty\}
\]
be its normalized additive valuation, defining the local-field topology:
$v(0)=+\infty$ and
$v(F^\times)=\Z$. Put
\[
 \OO=\{x\in F:v(x)\ge0\},\qquad
 \mm=\{x\in F:v(x)>0\},\qquad q=\#(\OO/\mm).
\]
Assume $q>2$. Fix $\varpi\in F^\times$ with $v(\varpi)=1$, and set
\[
 \abs a_F=q^{-v(a)}\quad(a\in F^\times).
\]
Let $\sqrt q$ denote the positive real square root. Let $\mu$ be the normalized
additive Haar measure on $F$ with $\mu(\OO)=1$, and let $\nu$ be the normalized
multiplicative Haar measure on $F^\times$ with
\[
 \nu(\OO^\times)=(1-q^{-1}).
\]
 All integrals below
use the product measure $\lambda=\nu\otimes(\mu\otimes\mu)$ on
$F^\times\times(F\times F)$.

Fix complex parameters $\alpha,\beta$ satisfying 
\begin{equation}\label{eq:regular}
 \abs\alpha=\abs\beta=1,\qquad
 \alpha\ne1,\quad\beta\ne1,\quad\alpha\ne\beta,\quad\alpha\beta\ne1.
\end{equation}


\clearpage
\section{The eight coefficients and the explicit kernel}

Define 
\begin{equation}\label{eq:C}
 C_q=1+2q^{-1}+2q^{-2}+2q^{-3}+q^{-4},\qquad
 R(z)=\frac{1-q^{-1}z}{1-z}.
\end{equation}
For each $0\le i\le7$, define complex numbers  $z_{i1},\ldots,z_{i4}$, $U_i,V_i$
by the following table, and put $A_i=\prod_{j=1}^{4}R(z_{ij})$.
\begin{center}
\renewcommand{\arraystretch}{1.32}
\begin{tabular}{@{}c c c c@{}}
\toprule
$i$ & $(z_{i1},z_{i2},z_{i3},z_{i4})$ & $U_i$ & $V_i$\\
\midrule
0 & $(\beta\alpha^{-1},\beta^{-1},\alpha^{-1}\beta^{-1},\alpha^{-1})$ & $\alpha\beta$ & $\alpha$\\
1 & $(\alpha\beta^{-1},\alpha^{-1},\alpha^{-1}\beta^{-1},\beta^{-1})$ & $\alpha\beta$ & $\beta$\\
2 & $(\alpha^{-1}\beta^{-1},\beta,\beta\alpha^{-1},\alpha^{-1})$ & $\alpha\beta^{-1}$ & $\alpha$\\
3 & $(\alpha\beta,\alpha^{-1},\beta\alpha^{-1},\beta)$ & $\alpha\beta^{-1}$ & $\beta^{-1}$\\
4 & $(\alpha^{-1}\beta^{-1},\alpha,\alpha\beta^{-1},\beta^{-1})$ & $\beta\alpha^{-1}$ & $\beta$\\
5 & $(\alpha\beta,\beta^{-1},\alpha\beta^{-1},\alpha)$ & $\beta\alpha^{-1}$ & $\alpha^{-1}$\\
6 & $(\beta\alpha^{-1},\alpha,\alpha\beta,\beta)$ & $\alpha^{-1}\beta^{-1}$ & $\beta^{-1}$\\
7 & $(\alpha\beta^{-1},\beta,\alpha\beta,\alpha)$ & $\alpha^{-1}\beta^{-1}$ & $\alpha^{-1}$\\
\bottomrule
\end{tabular}
\end{center}
The hypotheses \eqref{eq:regular} make every denominator in this table's
definition of $A_i$ nonzero.

For $a\in F^\times$ and $x,y,z\in F$, form the matrix
\begin{equation}\label{eq:matrix}
g= g(a;x,y,z)=
 \begin{pmatrix}
 1&0&y&z\\
 0&a&ax&ay\\
 0&0&a^{-1}&0\\
 0&0&0&1
 \end{pmatrix}.
\end{equation}
For such a matrix $g$, define
\begin{align}
 u(g)&=\min_{1\le i,j\le4}v(g_{ij}),\label{eq:u}\\
 s(g)&=\min_{\substack{1\le i_1<i_2\le4\\1\le j_1<j_2\le4}}
 v\bigl(g_{i_1j_1}g_{i_2j_2}-g_{i_1j_2}g_{i_2j_1}\bigr).\label{eq:s}
\end{align}
The minimum in \eqref{eq:s} is over all $36$ two-by-two minors. The valuation
of a zero entry or minor is $+\infty$.
Both displayed minima are finite integers. Put
\begin{equation}\label{eq:indices}
 n(g)=u(g)-s(g),\quad B(g)=s(g)-2u(g),\quad
 \ell(g)=2n(g),\quad H(g)=n(g)+B(g)=-u(g).
\end{equation}
Finally, \emph{define} the complex-valued function on these matrices by
\begin{equation}\label{eq:phi}
 \boxed{\displaystyle
 \Ph\bigl(g(a;x,y,z)\bigr)
 =\frac1{C_q}\sum_{i=0}^{7} A_i\,
 (q^{-1/2})^{6n(g)+4B(g)}U_i^{n(g)}V_i^{B(g)}.}
\end{equation}
The notation $\Ph$ in this document means precisely \eqref{eq:phi}.

\clearpage
\section{Integrands and the closed expression}

Set
\begin{equation}\label{eq:central}
 (z_0,z_1,z_2)=(0,\varpi^{-1},\varpi^{-2}),\qquad
 (c_0,c_1,c_2)=((1-q^{-1})^{-1},q,-q/(1-q^{-1})).
\end{equation}


For $\delta\in\C^\times$, a real number $T>0$, and $r\in\{0,1,2\}$, define the function $ f_{r,T,\delta}$ on $F^\times \times F \times F$ by
\begin{equation}\label{eq:integrand}
 f_{r,T,\delta}(a,x,y)
 =\delta^{v(a)}\bigl(q^{-v(a)}\bigr)^{3/2}
 T^{H(g(a;x,y,z_r))}\Ph\bigl(g(a;x,y,z_r)\bigr).
\end{equation}

Whenever the  function in \eqref{eq:integrand} is \emph{absolutely integrable}, put
\begin{equation}\label{eq:core}
 \I_T(\delta)=\sum_{r=0}^{2}c_r
 \int_{F^\times\times(F\times F)} f_{r,T,\delta}\,d\lambda.
\end{equation}
The theorem below proves absolute integrability on
every parameter range where it uses \eqref{eq:core}.

Let $\mathcal A=(\alpha,\alpha^{-1},\beta,\beta^{-1})$, regarded as an
ordered tuple so that products retain multiplicities. Define
\begin{align}
 N_1&=\prod_{\zeta\in\mathcal A}(1-\zeta q^{-1}),\label{eq:N1}\\
 N_2&=(1-\alpha\beta q^{-1})(1-\alpha^{-1}\beta q^{-1})
 (1-\alpha\beta^{-1}q^{-1})(1-\alpha^{-1}\beta^{-1}q^{-1}),\label{eq:N2}\\
 D_1(\delta)&=\prod_{\zeta\in\mathcal A}(1-\zeta\delta (q^{-1/2})),\label{eq:D1}\\
 D_2(\delta)&=\prod_{\zeta\in\mathcal A}(1-\zeta\delta^{-1}(q^{-1/2})).\label{eq:D2}
\end{align}
Whenever its denominator is nonzero, define the quotient
\begin{equation}\label{eq:E}
 \boxed{\displaystyle
 \E(\alpha,\beta,\delta)=
 \frac{\dfrac{(1+\alpha)^2(1+\beta)^2}{\alpha\beta}
       (1-q^{-1})N_1N_2}
      {C_qD_1(\delta)D_2(\delta)}.}
\end{equation}
All denominators used in the theorem are proved nonzero under its hypotheses.

Every limit denoted by $T\uparrow1$ below is a limit through
\emph{real} $T$ satisfying $0<T<1$. 

\clearpage
\section{The formalized theorem}

\begin{theorem}\label{thm:main}
For every $\delta\in\C$ and real $T$ with
\begin{equation}\label{eq:convergence}
 q^{-1/2}\le\abs\delta\le1,\qquad 0<T\le1,\qquad T<\abs\delta q^{1/2},
\end{equation}
all three functions $f_{r,T,\delta}$, for $r \in \{0,1,2\}$, are measurable and absolutely
integrable with respect to $\lambda$, and so $\I_T(\delta)$ is well defined and finite.  In particular, this holds for every
$0<T<1$ on $q^{-1/2}\le\abs\delta\le1$, and for $T=1$ on $q^{-1/2}<\abs\delta\le1$. 

Moreover, the following identities hold.
\begin{enumerate}
\item If $(q^{-1/2})<\abs\delta\le1$, then $C_qD_1(\delta)D_2(\delta)\ne0$ and
\begin{equation}\label{eq:principal}
 \boxed{\displaystyle
 \lim_{T\uparrow1}\I_T(\delta)=\I_1(\delta) = \E(\alpha,\beta,\delta).}
\end{equation}

\item Let $\varepsilon\in\{1,-1\}$ and assume additionally
\begin{equation}\label{eq:endpoint}
 \alpha\ne\varepsilon,\qquad\beta\ne\varepsilon.
\end{equation}
Then
$C_qD_1(\varepsilon (q^{-1/2}))D_2(\varepsilon (q^{-1/2}))\ne0$, and we have
\begin{equation}\label{eq:special}
 \boxed{\displaystyle
 \lim_{T\uparrow1}\frac{1-\varepsilon}{q+1}\I_T(\varepsilon (q^{-1/2}))
 =\frac{1-\varepsilon}{q+1}\E(\alpha,\beta,\varepsilon (q^{-1/2})).}
\end{equation}
\end{enumerate}
\end{theorem}

\begin{remark}

In part~(iv) of Theorem~\ref{thm:main}, the case $\varepsilon=1$
is trivial and simply states $0=0$.
In the case $\varepsilon=-1$, formula \eqref{eq:special} is the assertion
\[
 \lim_{T\uparrow1}\frac2{q+1}\I_T(-(q^{-1/2}))
 =\frac2{q+1}\E(\alpha,\beta,-(q^{-1/2})).
\]
The theorem does not assert the existence of $I_1(\delta)$ at either endpoint
$\delta=(q^{-1/2})$ or $\delta=-(q^{-1/2})$.
\end{remark}

\begin{remark}There are no additional analytic or representation-theoretic
hypotheses behind Theorem~\ref{thm:main}. All definitions and assumptions used are stated in this document.
\end{remark}


\end{document}
